\section*{TD 9 --- Synthèse : sujet de projet}
\addcontentsline{toc}{section}{TD 9 --- Synthèse : sujet de projet}

\subsection*{Exercice unique --- Cadrer un projet de bout en bout \hfill \normalfont\footnotesize (45 min)}
Une école veut prédire quels étudiants risquent d'échouer, pour proposer un
tutorat dès le premier mois. On dispose de cinq promotions, de la moyenne au
premier semestre, du nombre d'absences, de la filière et du résultat final.

\begin{enumerate}
  \item \textbf{Données.} Quelles variables sont disponibles au moment de la
        décision ? Lesquelles sont postérieures et doivent être écartées ?
  \item \textbf{Découpage.} Quel schéma de validation ? Justifier en une phrase.
  \item \textbf{Baseline.} Quelle baseline proposez-vous, et quel score
        obtiendrait-elle approximativement si 70 \% des étudiants réussissent ?
  \item \textbf{Métrique.} L'école veut éviter de rater un étudiant en
        difficulté. Quelle métrique privilégier, et quelles en sont les
        conséquences ?
  \item \textbf{Fuite.} Citer deux fuites possibles dans ce projet précis.
  \item \textbf{Usage.} Rédiger une phrase décrivant ce que l'école fera du
        score, et une phrase décrivant ce qu'elle ne doit pas en faire.
\end{enumerate}

\vspace{3mm}
{\footnotesize\color{gris}
Ce TD reprend la structure du projet intégrateur : cadrage, découpage,
baseline, métrique, fuite, usage. Un groupe qui répond correctement à ces six
questions a fait l'essentiel du travail de conception.}


\subsection*{Exercice 2 --- Coût et faisabilité \hfill \normalfont\footnotesize (10 min)}
\begin{enumerate}
  \item Estimer le volume de travail : collecte, annotation, modélisation,
        application.
  \item Quel est le risque principal du projet, et comment le réduire tôt ?
  \item Quel critère d'arrêt (go / no-go) poser au jalon 1 ?
\end{enumerate}

\subsection*{Exercice 3 --- Résultats attendus \hfill \normalfont\footnotesize (10 min)}
\begin{enumerate}
  \item Quel gain l'école peut-elle espérer, en points de rappel ?
  \item Quelle baseline faut-il battre pour que le projet ait un intérêt ?
  \item Comment mesurer que l'outil est réellement utilisé ?
\end{enumerate}

\subsection*{Exercice 4 --- Soutenance \hfill \normalfont\footnotesize (8 min)}
Préparer trois réponses courtes :
\begin{enumerate}
  \item « Pourquoi votre score est-il crédible ? »
  \item « Sur quels étudiants vous trompez-vous le plus ? »
  \item « Que fera l'école de vos prédictions ? »
\end{enumerate}

\newpage
